\chapter{Shin Splints (Medial Tibial Stress Syndrome)}

\section{Definition}
Shin splints (medial tibial stress syndrome) is exercise-induced pain along the shinbone (tibia). It is a common overuse injury in runners, dancers, and military recruits.

\section{Pathophysiology}
Repetitive stress on the tibia and its surrounding muscles causes periostitis (inflammation of the periosteum) at the insertion of the posterior tibialis and soleus muscles along the medial tibial border. As the stress continues, it can progress to a stress reaction and eventually a stress fracture of the tibia.

\section{Etiology}
\begin{itemize}
\item Sudden increase in exercise intensity, duration, or frequency
\item Running on hard or uneven surfaces
\item Improper footwear (worn-out shoes, inadequate support)
\item Overpronation of the feet (flat feet)
\item Tight calf muscles
\item Poor running mechanics
\item Inadequate warm-up
\item Female sex and low bone density (risk factors for progression to stress fracture)
\end{itemize}

\section{Indications for Evaluation}
\begin{itemize}
\item Dull, aching pain along the inner edge of the shinbone (tibia)
\item Pain during and after exercise
\item Pain that improves with rest
\item Tenderness to touch along the medial tibia
\item Mild swelling in the lower leg
\item Pain that persists at rest or at night (may indicate stress fracture)
\end{itemize}

\section{Related Symptoms}
\begin{itemize}
\item Stress fracture of the tibia
\item Compartment syndrome (acute or chronic exertional)
\item Tendonitis of the posterior tibialis
\item Calf muscle strain
\end{itemize}

\section{Self-Care/Home Management}
\begin{itemize}
\item Rest from high-impact activities
\item Apply ice packs for 20 minutes several times a day
\item Take OTC anti-inflammatory medications (ibuprofen, naproxen)
\item Switch to low-impact activities (swimming, cycling)
\item Stretch the calf muscles before and after activity
\item Massage the affected area gently
\item Consider orthotics for overpronation
\end{itemize}

\section{Diagnostic Approach}
\begin{itemize}
\item Clinical diagnosis based on history and tenderness along the medial tibia
\item X-ray may be normal early but can show stress fracture later
\item MRI for early detection of stress fracture or periosteal edema
\item Bone scan (scintigraphy) to differentiate shin splints from stress fracture
\end{itemize}

\section{Treatment/Management Overview}
\begin{itemize}
\item Relative rest: reduce intensity and frequency of activity
\item Ice and NSAIDs for pain and inflammation
\item Physical therapy focusing on stretching and strengthening
\item Proper footwear with good shock absorption and arch support
\item Custom orthotics for biomechanical issues
\item Gradual return to activity once pain-free
\item Cross-training with low-impact exercises
\end{itemize}
